\section{Discusion}
\label{sec:discussion}

\paragraph{Analisis de Sesgo Inductivo en Arquitecturas de Deep Learning:}

Los resultados obtenidos bajo condiciones controladas —todos los modelos entrenados desde cero con el mismo pipeline de normalizacion— revelan una jerarquia clara determinada por la presencia de sesgos inductivos convolucionales. CoAtNet (98.35\%), que combina etapas convolucionales MBConv con atencion Transformer, lidera el ranking al beneficiarse de las convoluciones en las primeras capas para extraer caracteristicas locales, mientras que la atencion en capas profundas modela dependencias de largo alcance. ResNet50 (97.48\%) y EfficientNet-B0 (97.00\%), arquitecturas puramente convolucionales, ocupan el segundo y tercer lugar respectivamente, demostrando que decadas de refinamiento en el diseno de CNNs producen modelos altamente eficientes incluso sin transfer learning.

En el extremo opuesto, ViT-B/16 (95.74\%) obtiene el peor desempeno a pesar de tener 85.8M de parametros —3.2 veces mas que CoAtNet y 21 veces mas que EfficientNet-B0—. La atencion global, que permite a cada parche relacionarse con todos los demas, se convierte en una desventaja cuando el volumen de datos es limitado: el modelo debe aprender desde cero que pixeles cercanos estan correlacionados (localidad) y que un mismo patron puede aparecer en cualquier posicion (equivarianza traslacional), conocimientos que las CNNs incorporan estructuralmente. Swin-T (96.32\%), con atencion restringida a ventanas locales de 7$\times$7, mitiga parcialmente esta deficiencia, superando a ViT-B/16 por 0.58 puntos porcentuales con 68\% menos parametros.

\paragraph{Implicaciones para el Diseno de Sistemas de Diagnostico Asistido:}

La diferencia de 2.61 puntos porcentuales entre CoAtNet (98.35\%) y ViT-B/16 (95.74\%) tiene implicaciones clinicas relevantes: en un conjunto de prueba de 1033 imagenes, CoAtNet cometio aproximadamente 17 errores frente a los 44 de ViT-B/16. Para un sistema de apoyo al diagnostico, esta reduccion de errores puede traducirse en menos falsos positivos (pacientes sanos derivados innecesariamente) y menos falsos negativos (tumores no detectados).

EfficientNet-B0 merece una mencion especial: con solo 4.0M de parametros y 55.3~ms de inferencia, alcanza 97.00\% de accuracy, superando a ViT-B/16 en 1.26 puntos porcentuales con el 4.7\% de sus parametros. Esto lo posiciona como la opcion ideal para despliegue en entornos con recursos limitados, donde la velocidad de inferencia y el consumo de memoria son criticos.

\paragraph{Limitaciones del Estudio:}

Si bien el pipeline de normalizacion elimino los artefactos estadisticos mas evidentes (diferencia de resolucion, modo de color y distribucion de brillo entre clases), la evaluacion se limito a un unico dataset de origen publico (Kaggle). La validacion externa con imagenes de otros centros hospitalarios o protocolos de adquisicion constituye un paso necesario para confirmar la generalizabilidad de los resultados. Asimismo, el entrenamiento desde cero con 30 epocas fijas representa una configuracion basal; la incorporacion de tecnicas adicionales como data augmentation mas agresiva o aprendizaje autosupervisado podria mejorar el rendimiento, aunque a costa de introducir nuevas variables de confusion en la comparacion.
